\section[H. Ordinary Offensive Magic]{\myboxviii{      \textbf{\texttt{H}}} Ordinary Offensive Magic}

\textit{Time limit}   2 seconds

\textit{Memory limit} 1 Gigabyte

\textit{Input} \qquad standard input

\textit{Output} \qquad standard output


\begin{quote}
	\textit{``Zoltraak.''}
\end{quote}

\subsection{Problem Statement}

Consider an $n \times m$ grid. Each cell has one of three states:

\begin{enumerate}
	\item \texttt{/} (forward-slash mirror): reflects on both sides. A ray turns $90^\circ$, connecting the top with the left and the bottom with the right.
	\item \texttt{\textbackslash} (backslash mirror): reflects on both sides. A ray turns $90^\circ$, connecting the top with the right and the bottom with the left.
	\item \texttt{.} (empty cell): contains no obstacle. A ray passes straight through, connecting the top with the bottom and the left with the right.
\end{enumerate}

Initially, some cells contain fixed \texttt{/} or \texttt{\textbackslash} mirrors; the rest are empty. You may \textbf{leave each initially empty cell empty}, \textbf{replace it with a \texttt{/} mirror}, or \textbf{replace it with a \texttt{\textbackslash} mirror}.

The outside boundary of the grid consists of $2(n+m)$ unit-length segments. A ray may enter or leave only perpendicularly through the midpoint of one of these segments.

Given an entry point and a target exit point, determine whether the grid can be filled legally so that a ray entering at the specified point, after straight-line travel and reflections, exits at exactly the target point.

If possible, output a valid completed grid that \textbf{minimizes the number of cells traversed} by the ray from the entry point to the target exit point. Count one traversal each time the ray enters a cell; repeated visits to the same cell are counted separately. If there are multiple optimal completions, output any of them. Otherwise, report that it is impossible.

\subsection{Input}

The first line contains an integer $T$ ($1 \le T \le 100$), the number of test cases.

The first line of each test case contains two integers $n,m$ ($1 \le n \times m \le 2 \times 10^5$), the numbers of rows and columns.

Each of the next $n$ lines contains a string of length $m$ describing the initial grid. Here \texttt{/} is a forward-slash mirror, \texttt{\textbackslash} is a backslash mirror, and \texttt{.} is an empty cell.

The next line contains four space-separated values $d_1,p_1,d_2,p_2$. The first pair specifies the entry point, and the second specifies the target exit point:

\begin{itemize}
	\item Each of $d_1,d_2$ is one of U, D, L, R, denoting the top, bottom, left, or right boundary, respectively.
	\item Each of $p_1,p_2$ is a positive integer specifying the midpoint of the outer edge of the $p$-th boundary cell. On U or D, $p$ is a column number ($1 \le p \le m$); on L or R, it is a row number ($1 \le p \le n$).
\end{itemize}

It is guaranteed that $(d_1,p_1) \ne (d_2,p_2)$.

The sum of $n \times m$ over all test cases does not exceed $2 \times 10^5$.

\subsection{Output}

For each test case, if a valid completion exists, output $n$ lines of length $m$ describing a completion that minimizes the number of cells traversed. If there are multiple optimal completions, output any of them. Do not output the minimum number of traversals. Otherwise, output \texttt{IMPOSSIBLE}.

A backslash must usually be escaped: \texttt{\textbackslash\textbackslash} outputs \texttt{\textbackslash}.

\subsection{Sample 1}

\begin{minipage}{\linewidth}
	\begin{minipage}[t]{0.49\linewidth}
		\begin{lstlisting}[title=\texttt{standard input}]
3
2 2
..
..
L 1 U 2
4 4
/\/\
/\/.
....
./..
L 2 L 4
4 4
/\/\
/\/.
..\/
./..
L 2 L 4
\end{lstlisting}
	\end{minipage}
	\begin{minipage}[t]{0.49\linewidth}
		\begin{lstlisting}[title=\texttt{standard output}]
./
..
/\/\
/\//
.//.
./..
IMPOSSIBLE
\end{lstlisting}
	\end{minipage}
\end{minipage}


\subsection{Sample 1 Explanation}

In the first test case, the ray enters from the left edge of row $1$ and travels right. It passes through the empty cell $(1,1)$, turns upward at the \texttt{/} mirror in $(1,2)$, and exits through the top edge of column $2$.

This completion traverses $2$ cells, which is the minimum possible.
